\documentclass[12pt,oneside]{amsart}

\usepackage{../qc-style}
\usepackage{../qc-macros}
\usepackage{../qc-tikz-templates}
\usepackage{listings}

% Julia code styling
\lstdefinelanguage{Julia}{
  keywords={function, end, for, while, if, else, elseif, return, using,
    import, export, module, begin, let, do, in, true, false, nothing},
  sensitive=true,
  comment=[l]{\#},
  morecomment=[s]{\#=}{=\#},
  string=[b]",
  morestring=[b]',
}
\lstset{
  language=Julia,
  basicstyle=\small\ttfamily,
  keywordstyle=\bfseries\color{cgblue},
  commentstyle=\itshape\color{spacecadet!60},
  stringstyle=\color{munsell},
  numbers=left,
  numberstyle=\tiny\color{spacecadet!40},
  numbersep=5pt,
  frame=single,
  framerule=0.4pt,
  rulecolor=\color{spacecadet!30},
  backgroundcolor=\color{isabelline},
  breaklines=true,
  columns=fullflexible,
  showstringspaces=false,
  xleftmargin=1.5em,
  xrightmargin=0.5em,
}

\title[Introduction to Quantum Computing --- Solutions]{Introduction to Quantum Computing\\[0.5em]
{\large Model Solutions}}
\author{Tobias J.\ Osborne}
\date{}

\begin{document}

\maketitle

%% ═══════════════════════════════════════════════════════════════
%%  PART I: QUANTUM CIRCUITS
%% ═══════════════════════════════════════════════════════════════
\section*{Part~I: Quantum circuits}
\setcounter{section}{1}
\setcounter{exercise}{0}

%% --- Exercise 1.1 ---
\begin{exercise}[Universality of $\NAND$]
\end{exercise}
\begin{solution}
  Recall $\NAND(a,b) = \overline{a \wedge b}$.  We simulate each gate:
  \begin{itemize}
    \item $\NOT(a) = \NAND(a,a) = \overline{a\wedge a} = \overline{a}$.
    \item $\AND(a,b) = \NOT(\NAND(a,b)) = \NAND(\NAND(a,b),\NAND(a,b))$.
    \item $\XOR(a,b) = \NAND(\NAND(a, \NAND(a,b)),\; \NAND(b, \NAND(a,b)))$.
  \end{itemize}
  For $\XOR$: let $c = \NAND(a,b)$, then $\NAND(a,c) = \overline{a\wedge\overline{ab}} = \overline{a}\vee (ab) = \overline{a}\vee b$ is not quite right in general, so let us verify by truth table.  Setting $c=\NAND(a,b)$, $d=\NAND(a,c)$, $e=\NAND(b,c)$, we have $\NAND(d,e)$:
  \begin{center}
    \begin{tabular}{cc|cccc}
      \toprule
      $a$ & $b$ & $c$ & $d$ & $e$ & $\NAND(d,e)$ \\
      \midrule
      0 & 0 & 1 & 1 & 1 & 0 \\
      0 & 1 & 1 & 1 & 0 & 1 \\
      1 & 0 & 1 & 0 & 1 & 1 \\
      1 & 1 & 0 & 1 & 1 & 0 \\
      \bottomrule
    \end{tabular}
  \end{center}
  This matches $\XOR(a,b) = a\oplus b$. \qedhere
\end{solution}

%% --- Exercise 1.2 ---
\begin{exercise}[$\CNOT$ circuit]
\end{exercise}
\begin{solution}
  In the computational basis $\{\ket{000},\ket{001},\ldots,\ket{111}\}$,
  $\CNOT_{12}$ flips qubit~2 when qubit~1 is $\ket{1}$, and
  $\CNOT_{23}$ flips qubit~3 when qubit~2 is $\ket{1}$.

  $\CNOT_{23}$ acts as $\I\tp \CNOT$ on qubits $(1,23)$:
  \[
    \CNOT_{23} = \begin{pmatrix}
      1&0&0&0&0&0&0&0\\
      0&1&0&0&0&0&0&0\\
      0&0&0&1&0&0&0&0\\
      0&0&1&0&0&0&0&0\\
      0&0&0&0&1&0&0&0\\
      0&0&0&0&0&1&0&0\\
      0&0&0&0&0&0&0&1\\
      0&0&0&0&0&0&1&0
    \end{pmatrix}.
  \]

  $\CNOT_{12}$ acts on qubits $(12,3)$ as $\CNOT\tp \I$:
  \[
    \CNOT_{12} = \begin{pmatrix}
      1&0&0&0&0&0&0&0\\
      0&1&0&0&0&0&0&0\\
      0&0&1&0&0&0&0&0\\
      0&0&0&1&0&0&0&0\\
      0&0&0&0&0&0&1&0\\
      0&0&0&0&0&0&0&1\\
      0&0&0&0&1&0&0&0\\
      0&0&0&0&0&1&0&0
    \end{pmatrix}.
  \]

  Thus $U = \CNOT_{12}\CNOT_{23}$ is the product of these matrices. The
  action on basis states is:
  $\ket{abc}\mapsto\ket{a,\;a\oplus b,\;(a\oplus b)\oplus c}$.

  The circuit diagram is:
  \begin{center}
  \begin{tikzpicture}
    \qcwires{3}{0}{6}
    \qcctrl{2}{1}
    \qccnottarg{2}{2}
    \qcctrlline{2}{1}{2}
    \qcctrl{4}{2}
    \qccnottarg{4}{3}
    \qcctrlline{4}{2}{3}
    \qclabel{1}{\ket{a}}
    \qclabel{2}{\ket{b}}
    \qclabel{3}{\ket{c}}
  \end{tikzpicture}
  \end{center}
\end{solution}

%% --- Exercise 1.3 ---
\begin{exercise}[Reversible simulation of classical circuits]
\end{exercise}
\begin{solution}
  Any classical gate $f:\{0,1\}^k\to\{0,1\}^l$ can be made reversible by
  using the \emph{Toffoli construction}: introduce $l$ ancilla bits
  initialised to $\ket{0}$ and compute
  $\ket{x}\ket{0^l}\mapsto\ket{x}\ket{f(x)}$.  Since
  $\NAND$ is universal for classical computation (Exercise~1.1), and the
  Toffoli gate (which is its own inverse) can simulate $\NAND$
  reversibly, any classical circuit can be simulated reversibly.

  Concretely, the Toffoli gate computes
  $\ket{a,b,c}\mapsto\ket{a,b,c\oplus(a\wedge b)}$.  With $c=1$ this gives
  $c\oplus(a\wedge b) = \NAND(a,b)$.  Composing reversible Toffoli gates to
  simulate any circuit of $\NAND$ gates gives a reversible simulation. The
  overhead is at most one ancilla bit per gate.
\end{solution}

%% --- Exercise 1.4 ---
\begin{exercise}[Controlled-$U$ matrix]
\end{exercise}
\begin{solution}
  Let $U = \begin{pmatrix} u_{00} & u_{01} \\ u_{10} & u_{11}\end{pmatrix}
  \in \SU{2}$.  The controlled-$U$ gate applies $U$ to the target qubit
  when the control is $\ket{1}$, and $\I$ when the control is $\ket{0}$:
  \[
    \CU{U} = \ketbra{0}{0}\tp \I + \ketbra{1}{1}\tp U
    = \begin{pmatrix}
      1 & 0 & 0 & 0 \\
      0 & 1 & 0 & 0 \\
      0 & 0 & u_{00} & u_{01} \\
      0 & 0 & u_{10} & u_{11}
    \end{pmatrix}.
    \qedhere
  \]
\end{solution}

%% --- Exercise 1.5 ---
\begin{exercise}[Two-level unitary decomposition]
\end{exercise}
\begin{solution}
  Let $U = \begin{pmatrix} a & d & g \\ b & e & h \\ c & f & j
  \end{pmatrix}$.

  \textbf{Step 1}: If $a \neq 0$, choose
  $U_1 = \begin{pmatrix} a^*/|a'| & b^*/|a'| & 0 \\ -b/|a'| & a/|a'| & 0
  \\ 0 & 0 & 1\end{pmatrix}$ where $|a'| = \sqrt{|a|^2+|b|^2}$.  Then
  $(U_1 U)_{21} = 0$.  If $a=0$, use a permutation to bring a non-zero
  entry into position $(1,1)$.

  \textbf{Step 2}: Apply a similar two-level unitary $U_2$ acting on rows
  $1$ and $3$ to zero out the $(3,1)$ entry of $U_2 U_1 U$.

  \textbf{Step 3}: After steps 1 and 2, the first column of $U_3' = U_2 U_1 U$
  is $(*, 0, 0)^T$.  Since $U_3'$ is unitary, it has the form
  $\begin{pmatrix} e^{i\alpha} & 0 \\ 0 & V \end{pmatrix}$ where $V$ is a
  $2\times 2$ unitary.  Set $U_3 = (U_3')^\dagger$, which is itself
  two-level.  Then $U_3 U_2 U_1 U = \I$.

  \textbf{Generalisation}: For $d\times d$, the first column has $d-1$
  entries to zero out (each requiring one two-level unitary), the second
  column has $d-2$, etc.  Total: $(d-1)+(d-2)+\cdots+1 = d(d-1)/2$.
\end{solution}

%% --- Exercise 1.6 ---
\begin{exercise}[Gray codes]
\end{exercise}
\begin{solution}
  We need a path from $\mathbf{s}=101001$ to $\mathbf{t}=110011$ changing
  one bit at a time.  The bits differ in positions 1, 2, 4, 6
  (numbering from left, starting at 1).  A valid Gray code:
  \[
    \begin{matrix}
      101001 \\
      111001 \\
      110001 \\
      110011
    \end{matrix}
  \]
  (Change bit~2, then bit~3, then bit~5---wait, let us recheck. Positions
  where $\mathbf{s}$ and $\mathbf{t}$ differ: position~1 ($1$ vs $1$---same),
  position~2 ($0$ vs $1$), position~3 ($1$ vs $0$), position~5 ($0$ vs $1$),
  position~6 ($1$ vs $1$---same).  Actually let us compare carefully:
  $\mathbf{s}=1\mathbf{0}1\mathbf{0}0\mathbf{1}$ vs
  $\mathbf{t}=1\mathbf{1}0\mathbf{0}1\mathbf{1}$.
  Differences at positions 2, 3, 5.  So:
  \[
    \begin{matrix}
      101001 \\
      111001 \\
      110001 \\
      110011
    \end{matrix}
  \]
  Change position~2: $101001\to 111001$. Change position~3: $111001\to 110001$.
  Change position~5: $110001\to 110011$.  Each step differs in exactly one
  bit. \qedhere
\end{solution}

%% --- Exercise 1.7 ---
\begin{exercise}[Implementing a two-level unitary]
\end{exercise}
\begin{solution}
  The matrix $U$ acts non-trivially on $\ket{000}$ and $\ket{110}$
  (positions 0 and 6 in the $8\times 8$ matrix). The Gray code
  $000\to 010\to 110$ provides the circuit construction.

  \textbf{Step 1}: Swap $\ket{000}\leftrightarrow\ket{010}$ by flipping
  qubit~2 conditioned on qubits 1 and 3 being $0$. This is a controlled-$X$
  on qubit~2, with controls on $\overline{q_1}$ and $\overline{q_3}$
  (i.e., controlled on qubits 1, 3 both being 0).

  \textbf{Step 2}: Apply controlled-$\widetilde{U}$ to qubit~1 (the bit
  where $010$ and $110$ differ), conditioned on qubit~2 being $1$ and
  qubit~3 being $0$.

  \textbf{Step 3}: Undo Step~1 (apply the same controlled-$X$ again, since
  it is self-inverse).

  The complete circuit uses two multi-controlled $X$ gates and one
  multi-controlled $\widetilde{U}$ gate, each decomposable into Toffoli and
  single-qubit gates.
\end{solution}

%% ═══════════════════════════════════════════════════════════════
%%  PART II: QUANTUM ALGORITHMS
%% ═══════════════════════════════════════════════════════════════
\section*{Part~II: Quantum algorithms}
\setcounter{section}{2}
\setcounter{exercise}{0}

%% --- Exercise 2.1 ---
\begin{exercise}[Phase oracle from standard oracle]
\end{exercise}
\begin{solution}
  We want to implement $\ket{j,b}\mapsto\ket{j,b\oplus x_j}$ using
  controlled phase queries $\ket{c,j}\mapsto(-1)^{c\cdot x_j}\ket{c,j}$.

  Prepare an ancilla in $\ket{+}=\frac{1}{\sqrt{2}}(\ket{0}+\ket{1})$.
  Apply the controlled phase query with this ancilla as the control:
  \[
    \ket{+}\ket{j} \xmapsto{\text{c-phase}}
    \frac{1}{\sqrt{2}}\bigl(\ket{0}\ket{j} + (-1)^{x_j}\ket{1}\ket{j}\bigr).
  \]
  Now apply a Hadamard to the ancilla:
  \[
    H\frac{1}{\sqrt{2}}\bigl(\ket{0}+(-1)^{x_j}\ket{1}\bigr)
    = \ket{x_j}.
  \]
  So the ancilla now holds $\ket{x_j}$.  Using a $\CNOT$ from this ancilla
  to a target qubit initialised in $\ket{b}$ gives $\ket{b\oplus x_j}$,
  implementing the standard query. Finally, uncompute the ancilla by
  reversing the steps.
\end{solution}

%% --- Exercise 2.2 ---
\begin{exercise}[Deutsch-Jozsa with workspace]
\end{exercise}
\begin{solution}
  \textbf{(a)} The circuit $H\Oxpm H$ on input $\kezero$:
  \[
    \kezero \xrightarrow{H} \frac{1}{\sqrt{2}}(\kezero+\keone)
    \xrightarrow{\Oxpm} \frac{1}{\sqrt{2}}((-1)^{x_0}\kezero+(-1)^{x_1}\keone)
    \xrightarrow{H} \frac{(-1)^{x_0}+(-1)^{x_1}}{2}\kezero
    + \frac{(-1)^{x_0}-(-1)^{x_1}}{2}\keone.
  \]
  The probability of measuring $0$ is
  $P(0) = \frac{1}{4}|(-1)^{x_0}+(-1)^{x_1}|^2$.
  \begin{itemize}
    \item If $x_0=x_1$ (constant): $P(0)=1$, $P(1)=0$.
    \item If $x_0\neq x_1$ (balanced): $P(0)=0$, $P(1)=1$.
  \end{itemize}
  This is exactly the Deutsch-Jozsa algorithm for $n=1$.

  \textbf{(b)} With workspace, the oracle acts as
  $O'_{x,\pm}\ket{b,0} = (-1)^{x_b}\ket{b,b}$.  Starting from
  $\ket{0,0}$:
  \begin{align}
    \ket{0,0} &\xrightarrow{H\tp \I}
    \frac{1}{\sqrt{2}}(\ket{0,0}+\ket{1,0}) \\
    &\xrightarrow{O'_{x,\pm}}
    \frac{1}{\sqrt{2}}((-1)^{x_0}\ket{0,0}+(-1)^{x_1}\ket{1,1}) \\
    &\xrightarrow{H\tp \I}
    \frac{(-1)^{x_0}}{2}(\ket{0}+\ket{1})\ket{0}
    +\frac{(-1)^{x_1}}{2}(\ket{0}-\ket{1})\ket{1}.
  \end{align}
  The probability of measuring the first qubit as $0$:
  \[
    P(0) = \left|\frac{(-1)^{x_0}}{2}\right|^2
    + \left|\frac{(-1)^{x_1}}{2}\right|^2 = \frac{1}{2}.
  \]
  Regardless of $x$, the output is uniformly random! The workspace qubit
  has \emph{decohered} the algorithm by entangling with the address
  register, destroying the interference needed for the algorithm to work.
\end{solution}

%% --- Exercise 2.3 ---
\begin{exercise}[Grover's algorithm step-by-step]
\end{exercise}
\begin{solution}
  With $n=2$, $N=4$, $x=0001$, the unique solution is $j=11$ (index~3).
  Thus $t=1$ and $\theta=\sin^{-1}(1/2)=\pi/6$.

  \textbf{Initial state}: $\ket{U} = \frac{1}{2}(\ket{00}+\ket{01}+\ket{10}+\ket{11})$.

  \textbf{$k=1$, Step (a) -- Oracle $\Oxpm$}: Flip the sign of $\ket{11}$:
  \[
    \frac{1}{2}(\ket{00}+\ket{01}+\ket{10}-\ket{11}).
  \]

  \textbf{$k=1$, Step (b) -- Diffusion $2\ket{U}\bra{U}-\I$}:
  The diffusion operator reflects about the mean.  The mean amplitude is
  $\bar{\alpha} = \frac{1}{4}(1/2+1/2+1/2-1/2) = 1/4$.
  Each amplitude is mapped to $2\bar{\alpha}-\alpha_j$:
  \[
    \frac{1}{2} \mapsto 2\cdot\frac{1}{4}-\frac{1}{2} = 0, \quad
    -\frac{1}{2} \mapsto 2\cdot\frac{1}{4}+\frac{1}{2} = 1.
  \]
  Final state after $k=1$: $\ket{11}$.

  \textbf{Success probability for $k=1$}: $P_1 = |\braket{11}{\psi_1}|^2 = 1$.

  This confirms $\widetilde{k}=1$ since $\theta=\pi/6$ gives
  $(2\cdot 1+1)\cdot\pi/6 = \pi/2$.

  \textbf{$k=2$}: Applying $\mathcal{G}$ again to $\ket{11}$:
  Oracle flips: $-\ket{11}$.  Diffusion about $\ket{U}$:
  mean $= -1/4$, so new amplitudes are $2(-1/4)-0=-1/2$ for non-solutions
  and $2(-1/4)-(-1) = 1/2$ for $\ket{11}$.
  State: $\frac{1}{2}(-\ket{00}-\ket{01}-\ket{10}+\ket{11})$.

  Success probability for $k=2$: $P_2 = |1/2|^2 = 1/4$.

  This illustrates the ``overshooting'' phenomenon: too many iterations move
  amplitude \emph{away} from the solution. The angle is
  $(2\cdot 2+1)\pi/6 = 5\pi/6$, giving $\sin^2(5\pi/6)=1/4$.
\end{solution}

%% --- Exercise 2.4 ---
\begin{exercise}[Grover with $t = N/4$]
\end{exercise}
\begin{solution}
  If $t = N/4$, then $\theta = \sin^{-1}(\sqrt{1/4}) = \sin^{-1}(1/2) = \pi/6$.
  After $k=1$ Grover iteration the state angle is
  $(2\cdot 1+1)\cdot\pi/6 = \pi/2$, so
  $P_1 = \sin^2(\pi/2) = 1$.  Thus one query suffices for certainty.

  \textbf{Classical with certainty}: In the worst case, a deterministic
  algorithm could query $3N/4$ non-solutions before finding a solution.
  So $3N/4+1 = 3N/4+1$ queries are needed in the worst case.
  More precisely, after $k$ queries all returning non-solutions, there remain
  $N-k$ unchecked positions, of which $N/4$ are solutions.  Certainty
  requires eliminating all non-solutions: $3N/4+1$ queries.

  \textbf{Classical with error $\le 1/10$}: We want
  $P(\text{fail after $k$ queries}) \le 1/10$.  If we query $k$ random
  positions, the probability of all being non-solutions is
  $\frac{\binom{3N/4}{k}}{\binom{N}{k}} \le (3/4)^k$.  We need
  $(3/4)^k \le 1/10$, giving $k \ge \frac{\log 10}{\log(4/3)} \approx 8$.
  So about $8$ random queries suffice.
\end{solution}

%% ═══════════════════════════════════════════════════════════════
%%  PART III: QUANTUM SIMULATION
%% ═══════════════════════════════════════════════════════════════
\section*{Part~III: Quantum simulation}
\setcounter{section}{3}
\setcounter{exercise}{0}

%% --- Exercise 3.1 ---
\begin{exercise}[Commuting Hamiltonians]
\end{exercise}
\begin{solution}
  \textbf{($\Rightarrow$)} Suppose $\comm{h_j}{h_k}=0$ for all $j,k$.
  Define $f(s) = \prod_{j=1}^L e^{-ih_j s}$ and
  $g(s) = e^{-iHs}$ for $s\in[0,t]$.  We show $f(s) = g(s)$ for all $s$.

  Differentiating: $f'(s) = \sum_{j=1}^L (-ih_j)\prod_{k=1}^L e^{-ih_k s}$.
  Since all terms commute, $h_j$ commutes with every $e^{-ih_k s}$, so
  $f'(s) = -i\sum_{j=1}^L h_j \cdot f(s) = -iH f(s)$.  Also $f(0) = \I$.
  Since $g(s) = e^{-iHs}$ satisfies the same ODE $g'(s) = -iH g(s)$ with
  $g(0)=\I$, uniqueness gives $f=g$.

  \textbf{($\Leftarrow$)} Conversely, suppose
  $e^{-iHt}=\prod_j e^{-ih_j t}$ for all $t$.  Expand both sides to second
  order in $t$:
  \begin{align}
    e^{-iHt} &= \I - iHt - \tfrac{1}{2}H^2 t^2 + \bigO{t^3}, \\
    \prod_j e^{-ih_j t} &= \I - iHt - \tfrac{1}{2}\sum_j h_j^2 t^2
    - \sum_{j<k} h_j h_k t^2 + \bigO{t^3}.
  \end{align}
  Comparing the $t^2$ coefficients:
  $H^2 = \left(\sum_j h_j\right)^2 = \sum_j h_j^2 + \sum_{j\neq k} h_j h_k$,
  while the product gives $\sum_j h_j^2 + 2\sum_{j<k} h_j h_k$.
  Equality requires $\sum_{j<k}(h_j h_k + h_k h_j) = 2\sum_{j<k} h_j h_k$,
  i.e., $\sum_{j<k}\comm{h_j}{h_k}=0$.  Since this must hold for all
  sub-collections, each commutator vanishes individually.
\end{solution}

%% --- Exercise 3.2 ---
\begin{exercise}[Counting local terms]
\end{exercise}
\begin{solution}
  If each $h_j$ acts on at most $c$ qubits, then $h_j$ is associated with a
  subset $S_j \subseteq \{1,\ldots,n\}$ with $|S_j|\le c$.  The number of
  distinct subsets of size at most $c$ is
  \[
    \sum_{k=0}^c \binom{n}{k} \le (c+1)\binom{n}{c}
    \le (c+1)\frac{n^c}{c!} = \bigO{n^c}.
  \]
  Since each term $h_j$ corresponds to one such subset (possibly with
  multiplicity bounded by $4^c$ due to the Pauli basis on $c$ qubits), the
  total number of terms is $L \le 4^c \binom{n}{c} = \bigO{n^c}$, which is
  polynomial in $n$ for fixed $c$.
\end{solution}

%% --- Exercise 3.3 ---
\begin{exercise}[Trotter error for the Ising model (computational)]
\end{exercise}
\begin{solution}
  See the Julia script \texttt{trotter\_ising.jl} for the complete
  implementation.  The key results are:

  \textbf{(a)} The exact unitary $U(t)$ is computed via matrix
  exponentiation of the full $16\times 16$ Hamiltonian matrix.

  \textbf{(b)} The Trotter approximation is computed by exponentiating each
  local term separately and taking the product, raised to the $n$th power.

  \textbf{(c)} The log-log plot of error vs.\ number of Trotter steps
  confirms the predicted $O(1/n)$ scaling (slope $\approx -1$ on the
  log-log plot).

  \begin{lstlisting}[caption={trotter\_ising.jl --- Trotter error analysis}]
using LinearAlgebra

# Pauli matrices
const I2 = [1.0+0im 0; 0 1]
const sx = [0.0+0im 1; 1 0]
const sz = [1.0+0im 0; 0 -1]

function kron_at(op, site, n)
    # Place operator `op` at `site` in n-qubit system
    foldl(kron, [i == site ? op : I2 for i in 1:n])
end

function ising_hamiltonian(n, lambda)
    d = 2^n
    H = zeros(ComplexF64, d, d)
    # XX coupling
    for j in 1:n-1
        H += kron_at(sx, j, n) * kron_at(sx, j+1, n)
    end
    # Transverse field
    for j in 1:n
        H += lambda * kron_at(sz, j, n)
    end
    return H
end

function ising_local_terms(n, lambda)
    terms = Matrix{ComplexF64}[]
    for j in 1:n-1
        push!(terms,
          kron_at(sx, j, n) * kron_at(sx, j+1, n))
    end
    for j in 1:n
        push!(terms, lambda * kron_at(sz, j, n))
    end
    return terms
end

function trotter_approx(terms, t, nsteps)
    dt = t / nsteps
    d = size(terms[1], 1)
    U_step = Matrix{ComplexF64}(I, d, d)
    for h in terms
        U_step = exp(-im * h * dt) * U_step
    end
    return U_step^nsteps
end

# Parameters
n_qubits = 4
lambda = 0.5
t = 1.0

H = ising_hamiltonian(n_qubits, lambda)
U_exact = exp(-im * H * t)
terms = ising_local_terms(n_qubits, lambda)

nsteps_list = [1, 2, 5, 10, 20, 50, 100]
for ns in nsteps_list
    U_approx = trotter_approx(terms, t, ns)
    err = opnorm(U_exact - U_approx)
    println("n=$ns, error=$(round(err, sigdigits=4))")
end
  \end{lstlisting}

  Representative output (operator norm error):
  \begin{center}
    \begin{tabular}{rc}
      \toprule
      $n$ (steps) & $\norm{U - \widetilde{U}}$ \\
      \midrule
      1   & $1.54$ \\
      2   & $0.796$ \\
      5   & $0.326$ \\
      10  & $0.164$ \\
      20  & $0.0824$ \\
      50  & $0.0330$ \\
      100 & $0.0165$ \\
      \bottomrule
    \end{tabular}
  \end{center}
  The ratio between successive errors confirms $O(1/n)$ scaling.
\end{solution}

\end{document}
